\section{Summary}
\label{sec:conclusions}

The Goldfarb--Idnani method is often presented as a sequence of updates, which
makes it look more intricate than it is. Read from the invariants it is short.

One matrix $J$ is carried, defined by $J\T G J = I$ and $J\T A = [R;0]$. The first
invariant says its columns are a $G$-orthonormal basis of $\R^n$; the second says
that basis is arranged so that the trailing columns span the working set's
feasible subspace and the leading ones the complement. Every quantity the
iteration needs is then one matrix--vector product or one triangular solve away,
and every quantity it \emph{consumes} --- the primal direction $z$, the dual
direction $r$ --- has a closed form in $(G, A, n_*)$ alone
(Proposition~\ref{prop:invariance}). That is why the representation may be chosen
on grounds of speed: two implementations holding different $J$ and $R$ for the same
working set compute the same step.

The iteration itself is one affine path (Proposition~\ref{prop:path}), along which
stationarity holds identically while the entering multiplier rises and the
working-set multipliers fall. Two limits cut it short, and which one binds decides
whether a constraint is added or dropped. The objective advances by an exact
closed form (Proposition~\ref{prop:increment}) which is positive in both step
directions, so the method is a dual ascent and, absent degeneracy, finite
(Corollary~\ref{cor:finite}). When the path is blocked in both directions the
vector $r$ is a Farkas certificate (Proposition~\ref{prop:farkas}): the
infeasibility verdict is a proof, and the solver already holds it.

Two properties of the classical method are easy to overlook and are what the
extensions of Sections~\ref{sec:pdas} and~\ref{sec:sweep} trade away or exploit.
The first is that linear dependence among the working-set normals is unreachable:
the step rules exclude it (Proposition~\ref{prop:rank}), so no rank test is
needed anywhere. Guessing the active set forfeits exactly that protection, which
is the real reason a primal--dual active-set fast path admits no
finite-termination theorem for general constraints --- there is no $P$-matrix
structure to appeal to once the working-set system stops being a principal
submatrix. The second is that the factors depend on $G$ and the working set but
not on the linear term, which is what makes a family of programs differing only in
$a$ solvable from one factorisation, and a stale active set repairable rather than
disposable.

Both extensions are unguaranteed and both are safe, for the same reason: strict
convexity makes the KKT conditions sufficient (Proposition~\ref{prop:sufficient}),
so a candidate can be certified rather than trusted. A method that may fail, but
that can recognise its own success, needs no convergence theory to be sound --- it
needs a fallback.

Three observations from Section~\ref{sec:implementation} generalise past this
solver, and all three are about the gap between a flop count and a running
program.

\emph{Layout decides which kernels are reachable.} Holding the triangular factor
packed rather than dense is worth a factor of $\qpLayRatio$ on the operation that
reads it (Section~\ref{ssec:packed}), and none of it is arithmetic. In a compiled
language layout is chosen for cache behaviour; in an array language it
additionally decides which library routine the data can be handed to at all, and a
layout that reaches only the general routine pays for that on every call.

\emph{A plausible mechanism is not a measured one.} The usual explanation for the
previous paragraph --- that a strided view forces a copy at the library boundary
--- is wrong here, and wrong in a way that casual measurement confirms rather than
exposes, because timing the two obvious shapes reproduces the headline ratio
(Remark~\ref{rem:notcopy}). The control that settles it costs one extra line and
we nearly did not run it.

\emph{Iterations and the cost of an iteration are different targets, and the
cheaper one is often the first.} Below a few hundred variables a solve costs what
its array operations cost to dispatch, so halving the arithmetic buys nothing and
halving the iteration count buys everything: the fast path of
Section~\ref{sec:pdas} beats the exact walk by more than the gap between two
implementations of the exact walk in different languages
(Section~\ref{ssec:compare}). Effort spent on kernels is wasted where the
bottleneck is the interpreter, and effort spent on iteration counts is wasted
where it is not; knowing which regime one is in is the whole of the decision.

Of the ideas here, the one that travels furthest is none of those. It is that a
method which may fail, but which can recognise its own success, needs a
certificate rather than a convergence theory.
